\section{Quality of Engineering Internships}\label{quality-of-engineering-internships}

\stanceinfo{\textbf{CFES Congress 2018} [2018-01-07]}{2024-01-02}

\officialstance{The Canadian Federation of Engineering Students believes that the quality of engineering internship programs across Canada is inconsistent and often substandard, and that engineering programs have a responsibility to revise their practices in order to offer students a better value for their program fees.}

\stanceheading{The Students' Position}
\begin{itemize}
 \item Student satisfaction with internship program quality is low across Canada, representing general dissatisfaction with the value of internship services offered at engineering schools including opportunities, resources provided and accessibility.
 \item The lack of regulation across engineering internship programs leads to greater variation and lower standards for the services that are offered to students, compared to other areas of engineering education.
 \item The flaws in internship programs vary considerably across institutions, meaning that an effective response to internship program issues must occur on a localized basis.
\end{itemize}

\stanceheading{Recommendations to Partners, Stakeholders, and Other Entities}
\begin{itemize}
 \item Student satisfaction with internship program quality is low across Canada, representing general dissatisfaction with the value of internship services offered at engineering schools including opportunities, resources provided and accessibility.
 \item The lack of regulation across engineering internship programs leads to greater variation and lower standards for the services that are offered to students, compared to other areas of engineering education.
 \item The flaws in internship programs vary considerably across institutions, meaning that an effective response to internship program issues must occur on a localized basis.
\end{itemize}

\stanceheading{What the CFES is doing}
\begin{itemize}
 \item The CFES included internship quality as a focus area for its 2017 National Student Survey and its associated research, in order to determine the scope of internship program issues across Canada.
 \item The CFES included internship and professional development questions in the 2023 National Survey to determine the requirements surrounding internship and resources available from institutions.
\end{itemize}

\stanceheading{What the CFES plans to do}
\begin{itemize}
 \item The CFES will work with its partners to create resources for sharing basic information about engineering program practices across Canada, and for advocating to faculties for effective changes.
 \item The CFES will investigate the merits of advocating to Co-operative Education and Work-Integrated Learning Canada to improve their accreditation of mandatory engineering internship programs.
 \item The CFES will bring back data collected by the National Survey regarding internships to Engineering Deans Canada in order to provide direction to their individual programs.
 \item The CFES will work with the regional organizations (WESST, ESSCO, QCESO, ACES) to gain information at a regional level of common engineering internship practice and policy.
\end{itemize}

\newpage
